
\section{Mathematical description of motion}

\subsection{Particles}

In this module we stick to the simplest models of Newtonian mechanics. In particular,
we approximate all bodies as point particles, that is, particles of zero size.
In doing so we are ignoring\\
~~~$\bullet$ that bodies can rotate\\
~~~$\bullet$ that bodies can have internal motions
\marginnote{
{\bf Reading:} Collinson and Roper \S 1.2. Note ideas of particle, speed, uniform motion,
mass, force.}

Often this is an excellent approximation (e.g.\ planets, rollercoaster). It is also an
important starting point for Newtonian mechanics because\\
$\bullet$ one can build up an {\it extended} body (one of non-zero size) by adding up
many point masses;\\
$\bullet$ given an extended body, there is a special point called the {\it centre of mass},
which moves as though all the mass were concentrated at that point.





\subsection{Position}

Given that we model a body as a point particle, we can define its {\bf position} at a
given time $t$. Introduce a set of fixed right-handed axes based at an origin~O.
\begin{marginfigure}
\includegraphics[width=5cm]{rhaxes.png}
\caption{\label{fig:rhaxes}
Right-handed axes.
}
\end{marginfigure}

At a given time the particle has position given by its coordinates $(x,y,z)$.
Since, in general, the particle moves, $x$, $y$ and $z$ vary with time, so we represent
position by three functions of time $(x(t),y(t),z(t))$.

We can equivalently represent the position of the particle at time $t$ using a
vector function of time $\boldsymbol{r}(t)$.
\begin{displaymath}
\boldsymbol{r}(t) = x(t) \boldsymbol{i} + y(t) \boldsymbol{j} + z(t) \boldsymbol{k} .
\end{displaymath}


\subsection{Velocity}

The {\bf velocity} of a particle is the rate of change of its position:
\begin{displaymath}
\boldsymbol{v}(t) = \frac{\dd \boldsymbol{r}}{\dd t} .
\end{displaymath}
At each time $t$, $\boldsymbol{v}$ is a vector whose magnitude gives the {\bf speed}
$|\boldsymbol{v} (t)|$ of the particle and whose direction gives the direction of
motion.
For brevity we will write $\boldsymbol{v} = \dot{\boldsymbol{r}}$; the dot indicates a
time derivative.
\begin{displaymath}
\boldsymbol{v} = \dot{\boldsymbol{r}} = \dot{x} \boldsymbol{i} + \dot{y} \boldsymbol{j} + \dot{z} \boldsymbol{k} .
\end{displaymath}

\ex{A particle has position}
\begin{displaymath}
\boldsymbol{r} = (\cos t^2 ) \boldsymbol{i} + (\sin t^2) \boldsymbol{j} .
\end{displaymath}
Find its velocity and speed as functions of time.


\subsection{Acceleration}

The {\bf acceleration} $\boldsymbol{a}(t)$ of a particle is the rate of change of its
velocity:
\begin{displaymath}
\boldsymbol{a}(t) = \frac{\dd \boldsymbol{v}}{\dd t}  = \frac{\dd^2 \boldsymbol{r}}{\dd t^2},
\end{displaymath}
or
\begin{displaymath}
\boldsymbol{a} = \dot{\boldsymbol{v}} = \ddot{\boldsymbol{r}} ,
\end{displaymath}
where two dots indicate two time derivatives.

\ex{Find the acceleration of the particle in the previous example.}


\subsection{Units}

In SI units, $x$, $y$, $z$, $\boldsymbol{r}$ are measured in metres~$\mathrm{m}$,\\
velocity and speed are measured in metres per second~$\mathrm{ms}^{-1}$,\\
and acceleration is measured in metres per second squared~$\mathrm{ms}^{-2}$.

%\noindent
%[{\bf Aside:}
\marginnote{
Mathematics questions don't always use units, but if a question uses units then
the answer should use units. Checking that the answer has the correct units can be a quick
way to identify that a mistake has been made.}

\ex{Particles $P_1$ and $P_2$ have position vectors
$\boldsymbol{r}_1 = u_0 t \boldsymbol{j} - g t^2 \boldsymbol{k} / 2$ and
$\boldsymbol{r}_2 = v_0 t \boldsymbol{i} - g t^2 \boldsymbol{k} / 2$, respectively,
where $u_0$, $v_0$, g are constants. Show that at $t=0$ the paths of the particles are
at right angles, while as $t \rightarrow \infty$ they become parallel.}

\vspace{10mm}

Note: in many problems we are free to choose where to place 
and how to orient our axes and when to
fix our origin in time. Good choices can often assist calculations.



\subsection{Validity of common sense notions of space and time}

We have appealed to `common sense' notions of time and space.
The physical world can be more subtle.

\vspace{3mm}

~~~~\begin{minipage}{100mm}
1. In relativity different observers disagree about measurements of
space and time (important when speeds approach the speed of light, or
when very large masses are involved).
\end{minipage}


\vspace{3mm}

~~~~\begin{minipage}{100mm}
2. In quantum mechanics particles do not have definite positions and velocities ---
these are replaced by probability distributions (important for very small systems
such as atoms, electrons,...).
\end{minipage}

\vspace{3mm}

\noindent
However, Newtonian mechanics and `common sense' notions are applicable to many
objects provided they are not moving too quickly (speed much less than that of light)
and not too small.


\marginnote{
\noindent
{\bf Reading:} Collinson and Roper \S 1.4 space and time, \S 4.1 motion in space.
}



\subsection{Initial conditions}

Given $\boldsymbol{r}(t)$, plainly we can find $\boldsymbol{v}(t)$ and $\boldsymbol{a}(t)$
by differentiating. Conversely, given $\boldsymbol{a}$, say, we can integrate to find
$\boldsymbol{v}$ and $\boldsymbol{r}$. This introduces (vector) constants of integration.
Often we can use {\bf initial conditions} or other information to fix these
constants.

\eg{A particle has acceleration
\begin{displaymath}
\boldsymbol{a} = t \boldsymbol{i} + t^2 \boldsymbol{j} + \alpha \boldsymbol{k}
\end{displaymath}
where $\alpha$ is a constant. Given that at $t=0$ the particle is at
the origin and has velocity $\boldsymbol{v} = \boldsymbol{i}$, obtain
the position and velocity at a general time.}

\ans{
First integrate $\boldsymbol{a}$:
\begin{eqnarray}
\boldsymbol{v} = \int \boldsymbol{a} \, \dd t
& = &
\int (t \boldsymbol{i} + t^2 \boldsymbol{j} + \alpha \boldsymbol{k}) \, \dd t
\nonumber \\
& = &
\frac{t^2}{2} \boldsymbol{i} + \frac{t^3}{3} \boldsymbol{j} + \alpha t \boldsymbol{k} + \boldsymbol{C} .
\nonumber
\end{eqnarray}
$\boldsymbol{C}$ is the constant of integration. Note that it is a \emph{vector} constant,
so written bold in print, underlined when hand-written.
\marginnote{
\emph{Always} underline vectors in your own working. Many unnecessary errors arise from
not taking enough care with notation, e.g.\ failing to distinguish
$v$, $\boldsymbol{v}$ and $V$, or $\omega$ and $w$.
}
At the moment it is unknown. But we do know
$\boldsymbol{v} = \boldsymbol{i}$ at $t=0$:
\begin{displaymath}
\boldsymbol{i} = 0 \boldsymbol{i} + 0 \boldsymbol{j} + 0 \boldsymbol{k} + \boldsymbol{C} .
\end{displaymath}
So $\boldsymbol{C} = \boldsymbol{i}$:
\begin{displaymath}
\boldsymbol{v} =
\left( 1 + \frac{t^2}{2} \right) \boldsymbol{i} + \frac{t^3}{3} \boldsymbol{j} + \alpha t \boldsymbol{k} .
\end{displaymath}
Integrate again
\begin{eqnarray}
\boldsymbol{r} = \int \boldsymbol{v} \, \dd t
& = &
\int \left\{
\left( 1 + \frac{t^2}{2} \right) \boldsymbol{i} + \frac{t^3}{3} \boldsymbol{j} + \alpha t \boldsymbol{k}
\right\} \, \dd t
\nonumber \\
& = &
\left( t + \frac{t^3}{6} \right) \boldsymbol{i} + \frac{t^4}{12} \boldsymbol{j} + \frac{\alpha t^2}{2} \boldsymbol{k}
+ \boldsymbol{D} \nonumber
\end{eqnarray}
where $\boldsymbol{D}$ is another (vector) constant of integration.
But $\boldsymbol{r} = 0$ at $t = 0$ $\Rightarrow$ $\boldsymbol{D} = \boldsymbol{0}$. Hence
\begin{displaymath}
\boldsymbol{r} =
\left( t + \frac{t^3}{6} \right) \boldsymbol{i} + \frac{t^4}{12} \boldsymbol{j} + \frac{\alpha t^2}{2} \boldsymbol{k} .
\end{displaymath}
}


\subsection*{Practice problems}

\vspace{3mm}
\refstepcounter{examplecounter}
\noindent \theexamplecounter
\label{Q:5.1}
A body has position vector 
$\boldsymbol{r}=\boldsymbol{i} \cosh t  + \boldsymbol{j} \cos (3t) + \boldsymbol{k} (t^2 +1)$.
Obtain its velocity, speed, and acceleration.

\vspace{3mm}
\refstepcounter{examplecounter}
\noindent \theexamplecounter
\label{Q:5.2}
Two particles have position vectors
$\boldsymbol{r}_1 (t) = \boldsymbol{i} \cos (2t) + \boldsymbol{j} \sin (2t) + \boldsymbol{k} ct$;
$\boldsymbol{r}_2 (t) = \boldsymbol{i} \cos (2t) - \boldsymbol{j} \sin (2t) + \boldsymbol{k} ct$.
At what time or times are the particle paths parallel?
What is the maximum value that $c$ may take while permitting the paths to become
perpendicular at some time?


%\newpage
\vspace{10mm}
\section{Newton's Laws}

So far we have not discussed what {\it causes} a body to move, accelerate, etc.
We now begin a study of {\bf Dynamics} (in its strictest sense---we will broaden
the topic later), which is concerned with forces acting on a body and how these
determine the motion.

In classical mechanics we begin with {\bf Newton's laws}:

\vspace{3mm}
\begin{minipage}{100mm}

\noindent
{\bf First law.} A body will remain at rest or continue to move in a straight line
with constant velocity unless a force acts on it.

\noindent
{\bf Second law.} When an external force is applied to a body, the force produces
an acceleration proportional to the force. The constant of proportionality is the
mass of the body.

\noindent
{\bf Third law.} When a body $A$ exerts a force on a body $B$ the body $B$ exerts an
equal and opposite force on $A$.

\end{minipage}
\vspace{3mm}

\noindent
(Note there are various alternative statements of Newton's laws. In particular in 
deeper treatments the first law becomes a statement about the existence of 
\textbf{inertial frames} in which particles move in straight lines when no forces act.)
Force is a vector, having magnitude and direction. It has units of newtons, symbol
$\mathrm{N}$. $1 \, \mathrm{N} = 1 \, \mathrm{kg}\,\mathrm{ms}^{-2}$.


\begin{figure*}
\hspace{3cm}Gravity \hspace{4cm}Electromagnetic forces

\includegraphics[width=60mm]{schematic_gravity.png}
~~~~~~~~
\includegraphics[width=60mm]{schematic_electromagnetic.png}
\vspace{0.5cm}

\hspace{3cm}Tension \hspace{4cm}Drag  or air resistance

\includegraphics[width=60mm]{schematic_spring.png}
~~~~~~~~~~
\includegraphics[width=60mm]{schematic_drag.png}
\vspace{0.5cm}

\hspace{4cm}Friction forces, normal reaction forces

~~~~~~~~~~~~~~~~~~~~~~~~~~
\includegraphics[width=80mm]{schematic_friction.png}
\caption{\label{fig:exampleforces}
Examples of forces.}
\end{figure*}


If there are several forces $\boldsymbol{F}_1$, $\boldsymbol{F}_2$, ... acting on a body
then the total force is
\begin{displaymath}
\boldsymbol{F} = \sum_i \boldsymbol{F}_i .
\end{displaymath}

Newton's second law relates the total force to the acceleration:
\begin{displaymath}
\mathrm{force}
\ = \ 
\mathrm{mass}
\ \times \ 
\mathrm{acceleration}
\end{displaymath}
i.e.\
\begin{displaymath}
\boldsymbol{F} =  m \boldsymbol{a} .
\end{displaymath}
Consequently, if $\boldsymbol{F} = \boldsymbol{0}$ (no external force acting) then
$\boldsymbol{a} = \boldsymbol{0}$, so, integrating,
\begin{displaymath}
\boldsymbol{v} = \int \boldsymbol{a} \, \dd t = \boldsymbol{c} = \ \mathrm{constant} .
\end{displaymath}
Thus Newton's first law follows from his second.
Mass is measured in kilograms ($\mathrm{kg}$). In Newtonian theory the mass
of a body is constant.

\eg{ A body of mass $3$ is subjected to a force $\boldsymbol{F} = \sin
  t \, \boldsymbol{i} + \ee^{-t} \, \boldsymbol{j}$.  Initially it is
  at rest at the origin. Find its subsequent motion.}

\vspace{5mm}
\ans{
By $\boldsymbol{F} = m \boldsymbol{a}$, the acceleration is
\begin{displaymath}
\boldsymbol{a} = \frac{1}{3} \sin t \, \boldsymbol{i} + \frac{1}{3} \ee^{-t} \, \boldsymbol{j} .
\end{displaymath}
Integrate:
\begin{displaymath}
\boldsymbol{v} = - \frac{1}{3} \cos t \, \boldsymbol{i} - \frac{1}{3} \ee^{-t} \, \boldsymbol{j} + \boldsymbol{C}.
\end{displaymath}
At $t = 0$, $\boldsymbol{v} = \boldsymbol{0}$, so
\begin{displaymath}
\boldsymbol{0} = - \frac{1}{3} \boldsymbol{i} - \frac{1}{3} \boldsymbol{j} + \boldsymbol{C} ,
\end{displaymath}
\begin{displaymath}
\boldsymbol{C} = \frac{1}{3} \boldsymbol{i} + \frac{1}{3} \boldsymbol{j} ,
\end{displaymath}
so
\begin{displaymath}
\boldsymbol{v} = \frac{1}{3} ( 1 - \cos t ) \boldsymbol{i} + \frac{1}{3} (1 - \ee^{-t} ) \boldsymbol{j} .
\end{displaymath}
Integrate again:
\begin{displaymath}
\boldsymbol{r} =
\frac{1}{3} ( t - \sin t ) \boldsymbol{i} + \frac{1}{3} (t + \ee^{-t} ) \boldsymbol{j} + \boldsymbol{D}.
\end{displaymath}
At $t = 0$, $\boldsymbol{r} = \boldsymbol{0}$, so
\begin{displaymath}
\boldsymbol{0} =
0 \boldsymbol{i} + \frac{1}{3} \boldsymbol{j} + \boldsymbol{D}
\ \ \ \ \Rightarrow \ \ \ \ 
\boldsymbol{D} = - \frac{1}{3} \boldsymbol{j} ;
\end{displaymath}
\begin{displaymath}
\boldsymbol{r} =
\frac{1}{3} ( t - \sin t ) \boldsymbol{i} + \frac{1}{3} (t - 1 + \ee^{-t} ) \boldsymbol{j} .
\end{displaymath}
}


\marginnote{{\bf Reading:} Collinson and Roper \S\S 4.2, 4.3.}


\subsection*{Practice problems}

\vspace{3mm}
\refstepcounter{examplecounter}
\noindent \theexamplecounter
\label{Q:6.1}
A body of mass $m$ is acted on by a force $\boldsymbol{F}(t)=\mu (1 - t) \boldsymbol{i}
+ \lambda t^2 \boldsymbol{j}$, where $\mu$ and $\lambda$ are constants;
obtain its motion if at $t=0$, $\boldsymbol{r}= X_0 \boldsymbol{i}$ and
$\boldsymbol{v}=V_0 \boldsymbol{k}$. 


\vspace{3mm}
\refstepcounter{examplecounter}
\noindent \theexamplecounter
\label{Q:6.2}
A small block of mass $m$  starts from rest on a slope of a hill
at an angle $\pi/6$ to the horizontal and at height $h$ above 
the horizontal. Find the acceleration of the block along the slope, 
and its speed along the slope when it reaches the foot of the hill. 
Assume there is a friction force of magnitude $mg/4$ pointing along the slope 
($g$ is the acceleration due to gravity).


\vspace{3mm}
\refstepcounter{examplecounter}
\noindent \theexamplecounter
\label{Q:6.3}
A particle of mass $m$ falls vertically under gravity (gravitational
acceleration $g$) and experiences a drag force due to air resistance
of magnitude $C |v|^2$ opposing the motion, where $v$
is the particle's speed and $C$ is a constant. Find the
terminal fall speed of the particle, i.e.\ the speed at which the
drag force balances the weight.


\section{Numerical estimates for derivatives}

We have seen how computers can help us to do calculations involving
(possibly lots of) arithmetic. But how can they help us with
problems involving calculus? Certain problems can be addressed
by symbolic manipulation packages such as Maple. Here we will learn
about another approach, in which a calculus problem is replaced
by a problem involving only arithmetic; specifically, derivatives are
replaced by arithmetic expressions that approximate the derivatives.
Variations on this idea underpin a great deal of scientific computing.

\subsection{First order Euler finite difference scheme}

By rearranging the Taylor series for $ f(x+h) $,
namely
%
\begin{equation}
f(x+h) = f(x) + h f'(x) + \tfrac{1}{2} h^2 f''(x) + \tfrac{1}{3!} h^3 f'''(x) + \cdots, \notag
\end{equation}
%
we obtain the Forward Euler formula
\begin{equation}
\label{eq:fwdeuler}
f'(x) = \frac{f(x+h) - f(x)}{h} + O(h).
\end{equation}
\marginnote{
The error is actually $\tfrac{1}{2} h f''(x)+ \cdots$  but the $O(h)$ notation for the 
error is convenient as it isolates the important information that the error
goes down proportional to $h$ when we reduce $h$.
}
Similarly, by rearranging the Taylor series for \( f(x-h) \), namely
%
\begin{equation}
f(x-h) = f(x) - h f'(x) + \tfrac{1}{2} h^2 f''(x) - \tfrac{1}{3!} h^3 f'''(x) + \cdots, \notag
\end{equation}
%
we obtain the Backward Euler formula
\begin{equation}
\label{eq:bwdeuler}
f'(x) = \frac{f(x) - f(x-h)}{h} + O(h).
\end{equation}
Both are first order accurate---i.e.\ the leading truncation error
scales like $ h $ to the power~1 as $ h \rightarrow 0 $.

\subsection{Second order centred difference scheme}

By combining the Taylor series for $ f(x+h) $ and $ f(x-h) $
we obtain the centred difference formula
\begin{equation}
\label{eq:cd2}
f'(x) = \frac{f(x+h) - f(x-h)}{2h} + O(h^2).
\end{equation}
This scheme is second order accurate.

\subsection{Fourth order centred difference scheme}

By combining the Taylor series for
$ f(x+2h) $, $ f(x+h) $, $ f(x-h) $, and $ f(x-2h) $,
we obtain the fourth order centred difference formula
\marginnote{
It is a good exercise to check this works and see
how the given coefficients ensure the cancellation of a key term.
}
%\vspace{-8mm}
\begin{displaymath}
f'(x) = \frac{-f(x+2h) + 8 f(x+h) -  8 f(x-h) + f(x-2h)}{12h} + O(h^4).
\end{displaymath}


\subsection{Centred difference formula for second derivative}

By combining Taylor series for $ f(x+h) $ and $ f(x-h) $
we can obtain a formula for the second derivative
\begin{equation}
\label{eq:d2dx}
f''(x) = \frac{f(x+h) - 2 f(x) + f(x-h)}{h^2} + O(h^2).
\end{equation}
This is second order accurate. 

\subsection{Convergence}

\marginnote{
\texttt{\# Examine convergence of}
}

\marginnote{
\texttt{\# Forward Euler}
}

\marginnote{
  \texttt{interval = []}
}

\marginnote{
  \texttt{err = []}
}

\marginnote{
\texttt{for p in range(8):}
}

\marginnote{
\texttt{\ \ \ \ h = 0.5**(p+1)}
}

\marginnote{
\texttt{\ \ \ \ d = (sin(1 + h) - sin(1))/h}
}

\marginnote{
\texttt{\ \ \ \ interval.append(h)}
}

\marginnote{
\texttt{\ \ \ \ err.append(abs(d - cos(1)))}
}

\marginnote{
\texttt{loglog(interval, err, 'b*')}
}


The order of accuracy of a finite difference scheme predicts its
convergence rate, as always, with the proviso that the function
is sufficiently smooth so that enough terms in the Taylor series
exist. For example, we can check the convergence of the Forward Euler
scheme on a function for which we know the true derivative.

\ex{Adapt this code to compare the convergence of Forward Euler
and centred differences. Do their convergence rates agree with their
orders of accuracy?}

\eg{

\marginnote{
\texttt{\# First load the data}
}
\marginnote{
\texttt{shtl = np.loadtxt('STS121ascent.txt')}
}
\marginnote{
\texttt{\# Pick out first column (time)}
}
\marginnote{
\texttt{time = shtl[:, 0]}
}
\marginnote{
\texttt{\# and second column (altitude)}
}
\marginnote{
\texttt{z = shtl[:, 1]}
}
\marginnote{
\texttt{\# How many data points?}
}
\marginnote{
\texttt{n = len(time)}
}
\marginnote{
\texttt{\# Estimate vertical velocity}
}
\marginnote{
\texttt{w = ...}
}
\noindent Download the shuttle data file from the ELE page and load
it into Python. Extract the first column (time) and second column (altitude)
of data. Use forward Euler to estimate the vertical velocity and
the vertical acceleration. Plot altitude, vertical velocity, and vertical
acceleration versus time.}



\begin{marginfigure}[1mm]
\begin{center}
\includegraphics[width=2.5cm]{KatherineJohnson.jpg}
\end{center}
\caption{\label{fig:KatherineJohnson}
Mathematician Katherine Johnson (1918--2020) performed complex manual
calculations and helped pioneer computer calculations
for NASA's space programmes.}
\end{marginfigure}


\newpage
\section{Gravity and projectiles}


\subsection{Newton's  law of universal gravitation}

Given two masses $m$ and $m'$, there is a force of attraction between them
proportional to $m m'$ and inversely proportional to the square of the distance
between them.


Suppose $m$ and $m'$ have position vectors $\boldsymbol{r}$ and $\boldsymbol{r}'$, respectively.
Then the force on $m$ due to $m'$ is
\begin{displaymath}
\boldsymbol{F} = - G m m' \frac{(\boldsymbol{r} - \boldsymbol{r}')}{|\boldsymbol{r} - \boldsymbol{r}'|^3} ,
\end{displaymath}
and the force on $m'$ due to $m$ is
\begin{displaymath}
\boldsymbol{F}' = - G m m' \frac{(\boldsymbol{r}' - \boldsymbol{r})}{|\boldsymbol{r} - \boldsymbol{r}'|^3} ,
\end{displaymath}
where $G$ is the gravitational; constant:
$G = 6.67 \times 10^{-11} \mathrm{m}^3\mathrm{s}^{-2}\mathrm{kg}^{-1}$.

\begin{marginfigure}[5mm]
\hspace{-0.9cm}
\includegraphics[width=75mm]{gravitymm.png}
\caption{\label{fig:gravitymm}
Newtion's law of gravitation.}
\end{marginfigure}

\vspace{5mm}
\noindent
Notes.

\noindent
1. $ |\boldsymbol{F}| = |\boldsymbol{F}'|
= \displaystyle{\frac{G m m'}{|\boldsymbol{r}' - \boldsymbol{r}|^2}}
= \displaystyle{\frac{G m m'}{( \mathrm{distance\ from}\ m\ \mathrm{to}\ m' )^2}}
$
\vspace{2mm}


\noindent
2. Since $m$, $m' > 0$, $\boldsymbol{F}$ and $\boldsymbol{F}'$ are in the directions shown;
the force is attractive.
\vspace{1mm}

\noindent
3. $\boldsymbol{F} = - \boldsymbol{F}'$, in agreement with Newton's third law.
\vspace{2mm}


When considering objects moving over moderate scales over Earth's surface, e.g.,
apples, balls, aeroplanes, shells, it is an excellent approximation to model the
gravitational force from the Earth on an object of mass $m$ as
\begin{displaymath}
\boldsymbol{F} = m \boldsymbol{g}
\end{displaymath}
where $\boldsymbol{g}$ is a constant vector pointing vertically downwards towards the
centre of the Earth. $|\boldsymbol{g}| = g \approx 9.81 \, \mathrm{ms}^{-2}$ is the
acceleration due to gravity.

We shall revisit the full inverse square law when we consider planetary motion;
however, for now we will take the gravitational force as $\boldsymbol{F} = m \boldsymbol{g}$.
This force is called the {\bf weight} of the body.


\subsection{Projectiles}
\label{sec:projectiles}

Consider a point particle (e.g.\ a ball or a shell), mass $m$, in motion, and suppose
the only force acting on it is gravity. We neglect air resistance for now.

Suppose the particle is launched (e.g.\ from a gun) from horizontal ground and
that its initial velocity is $\boldsymbol{V}$. To describe the motion we need to choose
some axes; a good choice simplifies matters.\\
(1) Take the time of launch to be $t=0$.\\
(2) Take the origin $O$ at the point of launch.\\
(3) Take $z$ to be vertically upwards with $z=0$ the ground, so that $\boldsymbol{g} = - g \boldsymbol{k}$.\\
(4) Take the $x$-axis such that $\boldsymbol{V}$ lies in the $(x,z)$-plane:
\begin{displaymath}
\boldsymbol{V} = V \cos \alpha \, \boldsymbol{i} + V \sin \alpha \, \boldsymbol{k} ,
\end{displaymath}
where $V = |\boldsymbol{V}|$ is the initial speed, and $\alpha$ is the angle of $\boldsymbol{V}$
to the horizontal.

Now, by Newton's second law $\boldsymbol{F} = m \boldsymbol{a}$
\begin{displaymath}
m \boldsymbol{a} = m \ddot{\boldsymbol{r}} = m \boldsymbol{g}
\ \ \ \Rightarrow \ \ \ 
\ddot{\boldsymbol{r}} = \boldsymbol{g} = \mathrm{constant}
\end{displaymath}
\begin{displaymath}
\boldsymbol{v} = \dot{\boldsymbol{r}} = \int \boldsymbol{g} \, dt = \boldsymbol{g} t + \boldsymbol{C}
\end{displaymath}
where $\boldsymbol{C}$ is the constant of integration.  At $t = 0$, $\dot{\boldsymbol{r}} = \boldsymbol{V}$
$\Rightarrow \ \ \boldsymbol{C} = \boldsymbol{V}$.
\begin{displaymath}
\dot{\boldsymbol{r}} = \boldsymbol{V} + \boldsymbol{g} t .
\end{displaymath}

Integrating again,
\begin{displaymath}
\boldsymbol{r} = \int ( \boldsymbol{V} + \boldsymbol{g} t ) \, dt
=
\boldsymbol{V} t + \frac{1}{2} \boldsymbol{g} t^2 + \boldsymbol{D} .
\end{displaymath}
At $t = 0$, $\boldsymbol{r} = \boldsymbol{0} \ \ \Rightarrow \ \ \boldsymbol{D} = \boldsymbol{0}$:
\begin{displaymath}
\boldsymbol{r} = \boldsymbol{V} t + \frac{1}{2} \boldsymbol{g} t^2 .
\end{displaymath}
This is the complete vector solution for the motion.

Now let us examine the solution in component form.
\begin{displaymath}
\begin{array}{llll}
\boldsymbol{V} & = V \cos \alpha \boldsymbol{i} &                & + V \sin \alpha \boldsymbol{k} \\
\boldsymbol{g} & =                          &                & - g \boldsymbol{k}             \\
\boldsymbol{r} & = x \boldsymbol{i}             & + y \boldsymbol{j} & + z \boldsymbol{k}             \\
\end{array}
\end{displaymath}
\vspace{2mm}
\begin{equation}
\boldsymbol{a} = \boldsymbol{g} \ \ 
\Rightarrow \ \ 
\left.
\begin{array}{lll}
\ddot{x} & = & 0 \\
\ddot{y} & = & 0 \\
\ddot{z} & = & -g \\
\end{array}
\right\}
\label{eq:accelxyz}
\end{equation}
\vspace{2mm}
\begin{displaymath}
\boldsymbol{v} = \dot{\boldsymbol{r}} = \boldsymbol{V} + \boldsymbol{g} t \ \ 
\Rightarrow \ \ 
\begin{array}{lll}
\dot{x} & = & V \cos \alpha \\
\dot{y} & = & 0             \\
\dot{z} & = & V \sin \alpha - g t \\
\end{array}
\end{displaymath}
\vspace{2mm}
\begin{displaymath}
\boldsymbol{r} = \boldsymbol{V} t + \frac{1}{2} \boldsymbol{g} t^2 \ \ 
\Rightarrow \ \ 
\begin{array}{lll}
x & = & V t \cos \alpha \\
y & = & 0               \\
z & = & V t \sin \alpha - \frac{1}{2} g t^2
\end{array}
\end{displaymath}

Note we could, instead, have started with \eqref{eq:accelxyz} and integrated, without
using vectors. Note also that $y=0$ for all times~$t$. This is a result of choosing the
initial velocity to be in the $(x,z)$-plane; once the motion starts in that plane it remains
there. Now that we have checked this, we can drop $y$ and just use $x$ and $z$ for
this type of problem.


\marginnote{
{\bf Reading:} Collinson and Roper \S 5.1 gravity, \S 4.4 projectiles.
}


The above solution
\begin{eqnarray}
x & = & V t \cos \alpha   \nonumber \\
z & = & V t \sin \alpha - \frac{1}{2} g t^2  \nonumber
\end{eqnarray}
gives the {\bf trajectory} of the particle as the curve $(x(t),z(t))$ parameterized by $t$.
To find $z = z(x)$ as a function of $x$, eliminate
\begin{equation}
t = \frac{x}{V \cos \alpha}
\label{eq:ttraj}
\end{equation}
to obtain
\begin{eqnarray}
z & = & x \tan \alpha - \frac{g x^2}{2 V^2 \cos^2 \alpha}  \nonumber \\
  & = & x \left( \tan \alpha - \frac{g x}{2 V^2 \cos^2 \alpha} \right) . \nonumber
\end{eqnarray}
This is the {\bf trajectory equation}.
$z$ is a quadratic function of $x$, so the curve is a {\bf parabola}. It passes through
the $x$-axis ($z=0$) at
\begin{displaymath}
x = 0 , \qquad \qquad x = \frac{2V^2}{g} \sin \alpha \cos \alpha = \frac{V^2}{g} \sin 2 \alpha .
\end{displaymath}
If the ground is horizontal then the {\bf range} of the projectile, i.e.\ the distance
before hitting the ground, is
\begin{displaymath}
R = \frac{V^2}{g} \sin 2 \alpha .
\end{displaymath}

The maximum value of $z$ along the trajectory, i.e.\ the {\bf height}~$h$, is
achieved when $\dot{z} = 0$, which occurs at $t = (V/g) \sin \alpha$,
and is given by
\begin{displaymath}
h = V t \sin \alpha - \frac{1}{2} g t^2 = \frac{V^2 \sin^2 \alpha }{2 g} .
\end{displaymath}

Also, from \eqref{eq:ttraj}, we know~$t$ as a function of~$x$, so when the projectile
lands at $x = R$ the {\bf time spent in flight} is
\begin{displaymath}
T = \frac{R}{V \cos \alpha} = \frac{2 V}{g} \sin \alpha .
\end{displaymath}

\vspace{5mm}

Now suppose we can vary $\alpha$, with $V$ and $g$ constant. Then the range~$R$ is
a maximum when $\sin 2 \alpha = 1$, which happens at $\alpha = \pi / 4$. In this
case $R = R_\mathrm{max} = V^2 / g$.

The height is a maximum when $\sin \alpha = 1 \ \ \Rightarrow \ \ \alpha = \pi/2$
and the particle is projected vertically upwards. Then $h = h_\mathrm{max} = V^2/(2g)$.

\eg{A particle is projected with velocity $V$ at an angle (a)~$\pi/6$,
  (b)~$\pi/3$ to the horizontal. Calculate its range, height, and time
  of flight in each case.}

\ans{
Note $\sin (\pi/3) = \sin (2\pi/3) = \sqrt{3}/2$, $\sin (\pi / 6) = 1/2$, so
for (a)
\begin{displaymath}
R = \frac{\sqrt{3}}{2} \frac{V^2}{g} , \qquad
h = \frac{V^2}{8 g} , \qquad
T = \frac{V}{g};
\end{displaymath}

\noindent
and for (b)
\begin{displaymath}
R = \frac{\sqrt{3}}{2} \frac{V^2}{g} , \qquad
h = \frac{3 V^2}{8 g} , \qquad
T = \frac{\sqrt{3} V}{g} .
\end{displaymath}
}

\begin{marginfigure}
\hspace{-0.8cm}
\includegraphics[width=65mm]{twotrajectories.png}
\caption{\label{fig:twotrajectories}
Two projectile trajectories.}
\end{marginfigure}

\eg{A ball is projected with speed $\sqrt{3gH}$ in a plane
  perpendicular to a vertical wall of height $H$ a distance $H$
  away. Find the range of angles of projection that allow the ball to
  pass over the wall.}

\ans{
Let $\alpha$ be the angle of projection, and we know the launch speed is $V = \sqrt{3 g H}$.
The equation of the trajectory is
\begin{eqnarray}
z & = & x \tan \alpha - \frac{g x^2}{2 V^2 \cos^2 \alpha} \nonumber \\
  & = & x \tan \alpha - \frac{x^2}{6 H} (1 + \tan^2 \alpha ) \nonumber
\end{eqnarray}
When $x = H$,
\begin{displaymath}
z = H \left( \tan \alpha - \frac{1}{6} - \frac{1}{6} \tan^2 \alpha \right)
\end{displaymath}
and for the ball to go over the wall this $z$ must be $\ge H$, i.e.\ we need
\begin{equation}
\tan \alpha - \frac{1}{6} - \frac{1}{6} \tan^2 \alpha \ge 1 .
\label{eq:trajwall}
\end{equation}
Let $Q = \tan \alpha$. Then we need $Q^2 - 6 Q + 7 \le 0$.
The roots of the quadratic are given by
\begin{displaymath}
Q = \frac{6 \pm \sqrt{36 - 28}}{2} = 3 \pm \sqrt{2} .
\end{displaymath}
Thus,~\eqref{eq:trajwall} is satisfied when $3 - \sqrt{2} \le \tan \alpha \le 3 + \sqrt{2}$.
The ball passes over the wall when
\begin{displaymath}
3 - \sqrt{2} \le \tan \alpha \le 3 + \sqrt{2}
\end{displaymath}
\begin{displaymath}
1.01 \le \alpha \le 1.35 .
\end{displaymath}
}

\subsection*{Practice problems}

\vspace{3mm}
\refstepcounter{examplecounter}
\noindent \theexamplecounter
\label{Q:7.1}
A long jumper attains a horizontal component of her velocity of
$8 \, \mathrm{ms^{-1}}$ at take-off. What is the minimum vertical component
of velocity she needs at take-off in order to jump a distance of $6 \, \mathrm{m}$ ?
In this case, what would her maximum height be?
[Take $g\simeq 10 \, \mathrm{ms}^{-2}$.] 


\vspace{3mm}
\refstepcounter{examplecounter}
\noindent \theexamplecounter
\label{Q:7.2}
A particle is projected with speed $3 (gh / 2)^{1/2}$ from a point~$O$ on a
horizontal floor. The particle reaches a point a horizontal distance~$3h$
from~$O$ on the edge of a table of height $h$. Find
the two possible angles of projection. 


\vspace{3mm}
\refstepcounter{examplecounter}
\noindent \theexamplecounter
\label{Q:7.3}
If a particle is projected from the horizontal floor of a large
room whose ceiling is at a height of $3 \, \mathrm{m}$, what will its greatest
possible range be if its speed of projection is (a)~$8 \,\mathrm{ms^{-1}}$,
(b)~$12 \, \mathrm{ms^{-1}}$ ?
[Take $g\simeq 10 \, \mathrm{ms}^{-2}$.] 


%\newpage
\subsection{The envelope equation}
\label{sec:envelope}

Suppose we have a gun that fires shells at a given speed $V$ and at any angle $\alpha$
to the horizontal. Which points of the $(x,z)$-plane can we hit?

We know that the trajectory is given by
\begin{displaymath}
z = x \tan \alpha - \frac{g x^2}{2 V^2 \cos^2 \alpha} .
\end{displaymath}
Thus it goes through a point $(X,Z)$ if
\begin{displaymath}
Z = X \tan \alpha - \frac{g X^2}{2 V^2 \cos^2 \alpha} ,
\end{displaymath}
i.e.\ if
\begin{displaymath}
Z = X \tan \alpha - \frac{g X^2}{2 V^2} (1 + \tan^2 \alpha),
\end{displaymath}
i.e.\ if
\begin{displaymath}
g X^2 \tan^2 \alpha - 2 V^2 X \tan \alpha + (2 V^2 Z + g X^2) = 0 .
\end{displaymath}
Given $g$ and $V$, we can hit $(X,Z)$ if this equation has a real solution for $\alpha$.
The equation is a quadratic in $\tan \alpha$ and has solutions.
\begin{eqnarray}
\tan \alpha & = & \frac{-B \pm (B^2 - 4 A C)^{1/2}}{2 A}  \nonumber \\
            & = & \frac{V^2 X \pm (V^4 X^2 - g X^2 (2 Z V^2 + g X^2))^{1/2}}{g X^2} .
\end{eqnarray}

There are three possibilities
\begin{enumerate}
\item
The quadratic has two real roots: there are two possible angles for which the trajectory
passes through $(X,Z)$.

\item
The quadratic has one repeated root: just one value of $\alpha$ allows us to hit $(X,Z)$.

\item
The quadratic as no real roots: it is impossible to hit $(X,Z)$ for any $\alpha$.
\end{enumerate}
Which case occurs depends on the discriminant $\Delta = B^2 - 4 A C$ of the quadratic.
If $\Delta / 4 = V^4 X^2 - g X^2 (2 Z V^2 + g X^2)$ \\
... is positive then we have two real roots, case 1; \\
... is zero then we have one real root, case 2; \\
... is negative then we have no real roots, case 3.\\
The key case is case 2, which occurs when
\begin{displaymath}
V^4 X^2 = g X^2 (2 Z V^2 + g X^2) ,
\end{displaymath}
i.e.,
\begin{equation}
\label{eq:envelope}
Z = \frac{V^2}{2 g} - \frac{g X^2}{2 V^2} .
\end{equation}
All points that can be hit with just one angle~$\alpha$ lie on this curve~$\mathcal{C}$.
If we sketch the curve we have\\
$X = 0 \ \ \Rightarrow \ \ Z = V^2 / 2 g$ (maximum height for gun); \\
$Z = 0 \ \ \Rightarrow \ \ X = \pm V^2 / g$ (maximum range for gun).
\begin{figure}
\hspace{-2cm}
\includegraphics[width=170mm]{envelope.png}
\caption{\label{fig:envelope}
The envelope of trajectories.}
\end{figure}
The curve $\mathcal{C}$ is called the {\bf envelope} of the trajectories.
It divides the $(x,z)$-plane into the region within range and the region out
of range. Equation~\eqref{eq:envelope} is the {\bf envelope equation}.

\eg{A ball is projected with speed $V$ in a plane perpendicular to a
  wall of height $H$ a distance $H$ away. What is the {\it minimum}
  speed required for the ball to pass over the wall?}

\ans{Given $V$, we can obtain the envelope~$\mathcal{C}$ of the
  trajectories launched from~$O$ with this speed $V$.
\begin{marginfigure}
\hspace{-1.cm}
\includegraphics[width=70mm]{envelope_over.png}

\hspace{-1.3cm}
\includegraphics[width=73mm]{envelope_under.png}
\caption{\label{fig:evelope_over_under}
The evelope may lie above or below the top of the wall.}
\end{marginfigure}
For the ball to go over the wall, the top of the wall $(X,Z) = (H,H)$ must lie
inside or on~$\mathcal{C}$. At the minimum possible speed $V$, the top of the wall
$(X,Z) = (H,H)$ must lie {\it on}~$\mathcal{C}$. Hence
\begin{displaymath}
H = \frac{V^2}{2 g} - \frac{g H^2}{2 V^2} ,
\end{displaymath}
\begin{displaymath}
2 = \frac{V^2}{gH} - \frac{g H}{V^2} .
\end{displaymath}
Substitute $\mu = V^2 / (g H)$, and solve the resulting quadratic equation in $\mu$
to obtain
\begin{displaymath}
\mu = 1 \pm \sqrt{2} .
\end{displaymath}
Since $V^2$, $g$ and $H$ are all positive, only the positive root makes physical sense,
i.e., $\mu = 1 + \sqrt{2}$. Therefore the minimum speed is
\begin{displaymath}
V = \sqrt{\mu g H} = \left( (1 + \sqrt{2}) g H \right)^{1/2} .
\end{displaymath} 
}

\subsection*{Practice problems}

\vspace{3mm}
\refstepcounter{examplecounter}
\noindent \theexamplecounter
\label{Q:7.4}
A particle is launched with speed $V = \sqrt{4gh}$ in a plane perpendicular
to a wall located a horizontal distance $h$ from the point of launch.
What is the maximum height of the wall over which the particle can pass?


\vspace{3mm}
\refstepcounter{examplecounter}
\noindent \theexamplecounter
\label{Q:7.5}
A gun is located at $x = 0$ on top of a hill whose height $z(x)$ is
described by $z = 4h - x^2 / h $, where $h$ is a constant. The gun fires shells
at speed $V = \sqrt{g h} / 2$. (i)~Determine the value of $x$ for the
most distant point on the hillside that the gun can hit.
(ii)~A target is located on the hillside at~$x = 2h$. Determine the 
smallest value of~$x$
to which the gun must be moved so that the target is just within range.


\section{Projectiles with air resistance: part~I}
\label{sec:air_res_I}

Real projectiles are subject to air resistance. A reasonable mathematical model
for the aerodynamic drag force on a particle is
\begin{equation}
\boldsymbol{F}_\mathrm{D} = - C |\boldsymbol{v}| \boldsymbol{v} ,
\end{equation}
where $\boldsymbol{v}$ is the velocity of the particle and $C = \rho_\mathrm{a} C_\mathrm{D} A$,
with $\rho_\mathrm{a}$ the air density, $C_\mathrm{D}$ a drag coefficient, and $A$ the
cross-sectional area of the particle\footnote{You may notice an apparent inconsistency
between treating a body as a point particle and allowing it to have a finite cross-sectional
area. In this case it does not lead to any problems. But, more generally, we do need to be
on the lookout for possible inconsistent assumptions and approximations in our modelling.
}.
For our purposes we can take $\rho_\mathrm{a}$, $C_\mathrm{D}$, and $A$, and hence $C$, as
constants. Newton's second law then becomes
\begin{equation}
\label{eq:proj_air_res}
m \ddot{\boldsymbol{r}} = - m \boldsymbol{g} - C |\boldsymbol{v}| \boldsymbol{v} .
\end{equation}
Although strictly the \emph{drag coefficient} is $C_{\mathrm{D}}$ as defined above, it is 
convenient in the informal discussion below to refer to $C$ as the drag coefficient.  Note that
\begin{equation}
|\boldsymbol{F}_\mathrm{D} |= - C |\boldsymbol{v}|^2
\end{equation}
and so the magnitude of the drag is quadratic in the speed of the body.

\marginnote{
Calculation of the drag  on a body is a topic in \emph{fluid dynamics}. Quadratic drag is appropriate 
for bodies such as cricket balls moving through the air (and we are assuming no wind), but Magnus effects 
can also arise from rotation. For a body moving in a viscous fluid, for example a bead sinking in liquid honey, 
the force is linear $|\boldsymbol{F}_\mathrm{D} | \propto |\boldsymbol{v}|$.
}

We can find closed form solutions to this equation only in three special cases:
\begin{enumerate}
\item
when $C = 0$, as we have seen above;
\item
when $\boldsymbol{g} = \boldsymbol{0}$;
\item
when the motion is purely vertical.
\end{enumerate}

When $\boldsymbol{g} = \boldsymbol{0}$ the direction of the particle's velocity never
changes and the motion is effectively one-dimensional. If we choose axes so that
the motion is in the positive $x$-direction then the equation of motion becomes
\begin{displaymath}
m \ddot{x} = - C \dot{x}^2 ,
\end{displaymath}
which can be integrated twice to obtain
\begin{displaymath}
x = \frac{m}{C} \log \left( 1 + \frac{t}{t_0} \right) ,
\end{displaymath}
where $t_0$ is a constant of integration related to the initial speed, and we have
assumed $x=0$ at $t=0$. [Exercise: check that this solution satisfies the equation of
motion.] 

When the motion is purely vertical then again we have a one-dimensional problem:
\begin{equation}
\label{eq:vertical_with_drag}
m \ddot{z} = - m g \pm C \dot{z}^2 ,
\end{equation}
with a minus sign when the motion is upward and a plus sign when the motion is downward.
It turns out that this, too, can be integrated to give
\begin{displaymath}
z = \left\{
\begin{array}{ll}
\ \ \ \frac{m}{C} \log \left\{ \cos \left( \sqrt{gC/m} \, t \right) \right\}  &
- \pi / 2 \sqrt{gC/m} < t \le 0 ,\\
- \frac{m}{C} \log \left\{ \cosh \left( \sqrt{gC/m} \, t \right) \right\}  &
t > 0 ,
\end{array}
\right.
\end{displaymath}
where the constants of integration have been chosen so that the particle comes to rest
at $z=0$ at $t=0$. For the first part the particle is moving upwards and $t<0$, then for 
$t>0$ it is moving downwards, while  at $t=0$ it is at maximum height.
[Exercise: check that this solution satisfies the equation of motion---take care to
choose the appropriate sign in \eqref{eq:vertical_with_drag}.]

For more general cases our only option is to find a numerical solution.
So we now need to look at numerical methods for initial value problems.


\section{Numerical methods for Initial Value Problems}
\label{sec:time_stepping}

The methods we look at in this section are often referred to as
\textit{time stepping} or \textit{time integration} methods.
To begin with we consider a first order (i.e.\ a single derivative)
ordinary differential equation (ODE) for a single scalar variable
\begin{equation}
\label{eq:ODE1}
\frac{d y}{d t} = f(y,t) .
\end{equation}
We will see below that the ideas extend straightforwardly to systems of
equations (section~\ref{sec:step_systems}) and to higher order
equations (section~\ref{sec:step_higher_systems}).

Let us try to find a numerical solution at a set of uniformly spaced times
$\{ t_0,\ t_1,\ \ldots,\ t_n,\ \ldots \}$, where $t_n = nh$ and $h$ is the
time step. Let $y_n$ be our approximate numerical solution at time $t_n$.
We assume that we are given the initial condition $y_0$.


\subsection{Single-step finite difference methods}

Now we need some way to calculate $y_{n+1}$ from $y_n$; this will allow us
to step forward from $y_0$ to $y_1$, from $y_1$ to $y_2$, and so on until
we have built up a solution over the desired time range.


\subsubsection{Forward Euler}

\marginnote{
\texttt{\# Solve dy/dt = y(1 - y)}
}

\marginnote{
\texttt{\# using Forward Euler}
}

\marginnote{
\texttt{\# Final time and time step}
}

\marginnote{
\texttt{Tstop = 5}
}

\marginnote{
\texttt{nstop = 50}
}

\marginnote{
\texttt{h = Tstop/nstop}
}

\marginnote{
\texttt{\# Initial value}
}

\marginnote{
\texttt{y = 0.1}
}

\marginnote{
\texttt{\# Time step loop}
}

\marginnote{
\texttt{for n in range(nstop):}
}

\marginnote{
\texttt{\ \ \ \ y = y + h*y*(1 - y)}
}

The simplest time stepping method is the forward Euler method. 
It is obtained by using the forward Euler approximation \eqref{eq:fwdeuler}
for the time derivative
\begin{displaymath}
\frac{\dd y}{\dd t} (t_n) \approx \frac{y_{n+1} - y_n}{h}.
\end{displaymath}


\noindent
The ODE \eqref{eq:ODE1} then becomes
\begin{equation}
y_{n+1} = y_n + h f(y_n, t_n) .
\label{eq:fwdeulerstep}
\end{equation}
Thus, if we know $y_n$ then we can evaluate the right hand side and hence
obtain $y_{n+1}$.
Starting with the initial condition $y_0$, we can successively find $y_1$, $y_2$, etc.
This idea is at the heart of all
single-step finite difference methods. 

\ex{Use the Forward Euler method to obtain a numerical solution of
$dy/dt = y(1-y)$ at $t = 5$ given $y = 0.1$ at $t = 0$. Take $h = 0.1$.}

\subsubsection{Midpoint method}

There are many other single-step methods.
\marginnote{For example, look up Runge-Kutta methods.}
Here we mention just one other: the midpoint method.

Here we attempt to calculate $dy/dt$ in the middle of the timestep, at $t=t_n+h/2$.
If the present moment is $t_n$, then $t=t_n+h/2$ is in the future. So first we need an
estimate of $y$ at the future time $t=t_n+h/2$, $y = y_*$, say. We can use
Forward Euler to obtain this estimate. Thus, the formula is
\begin{eqnarray*}
	y_* & = & y_n + \frac{h}{2} f(y_n, t_n),\\
	y_{n+1} & = & y_n + h f\left(y_*, t_n+ h/2 \right).
\end{eqnarray*}


\subsection{Multi-step methods}

If we use the second order centred difference method \eqref{eq:cd2} to approximate the
time derivative then the ODE \eqref{eq:ODE1} becomes
\begin{equation}
y_{n+1} = y_{n-1} + 2hf(y_n,t_n) .
\end{equation}
This method is known as the leapfrog method. It requires us to know both $y_{n-1}$
and $y_n$ in order to step forward to $y_{n+1}$, so it is an example of a
\textit{multi-step method}. 
\marginnote{Other multi-step schemes include the Adams-Bashforth methods.}
Note that in order to apply the leapfrog scheme we must do something else,
such as Forward Euler, for the first time step, since we don't have $y_{-1}$.


\subsection{Orders of accuracy and convergence}

Using the Taylor series \eqref{eq:taylor_series} shows that the error made in
one Forward Euler time step \eqref{eq:fwdeulerstep} is
$O(h^2)$. However, the number of steps needed to reach a given time~$T$
is $\propto 1/h$, and the errors will accumulate step by step.
Therefore the error in the solution at time $T$ is 
expected to be $O(h)$. Thus, the Forward Euler method is said to be first-order
accurate.
This analysis predicts that, if we repeat the calculation using different values of $h$,
the error at time~$T$ should get smaller, i.e.\ the numerical solution should {\it converge}
to the true solution, at a certain rate as $h \rightarrow 0$.

For both the midpoint method and the leapfrog method, Taylor series analysis
shows that the methods are second-order accurate. Thus, the error at time~$T$
is expected to be $O(h^2)$, etc.

Note that these predictions depend on the validity of the Taylor series expansion
from which the schemes are derived. Worse convergence may be achieved in practice
if the solution is not sufficiently smooth.
The predictions can also fail if the scheme is unstable---see section~\ref{sec:stability}.


\subsection{Time stepping systems of equations}
\label{sec:step_systems}

To apply the above time stepping methods to systems of equations we simply
replace the scalar variable $y$ and function $f$ by arrays or vectors
$\boldsymbol{y}$ and $\boldsymbol{f}$:
\begin{equation}
\label{eq:ODE12}
\frac{d \boldsymbol{y}}{d t} = \boldsymbol{f}(\boldsymbol{y},t) .
\end{equation}
For example, the forward Euler method becomes
\begin{equation}
\boldsymbol{y}_{n+1} = \boldsymbol{y}_n + h \boldsymbol{f}(\boldsymbol{y}_n, t_n) .
\label{eq:fwdeulerstepsys}
\end{equation}
For coding it may be convenient to split the equations into their components.

\eg{Find and plot a numerical solution of the SIR model over 200~days.
  Use the values of $S_0$, $I_0$, $r$, and $a$ given in
  section~\ref{sec:built-in}.}

\vspace{5mm}
\begin{minipage}{7cm}
\begin{verbatim}
# Solve SIR using Forward Euler

import matplotlib.pyplot as plt

# Final time and time step
Tstop = 200
nstop = 200
h = Tstop/nstop

# Constants
r = 3e-9
a = 1/14

# Initial values
t = [0]
S = [6e7]
I = [1000]
R = [0]
\end{verbatim}
\end{minipage}
\begin{minipage}{8cm}
\begin{verbatim}
# Time step loop
for n in range(nstop):
    t.append(t[n] + h)
    S.append(S[n] - h*r*I[n]*S[n])
    I.append(I[n] + h*(r*I[n]*S[n] - a*I[n]))
    R.append(R[n] + h*a*I[n])

# Plot results
plt.plot(t, S, 'b', label='S')
plt.plot(t, I, 'r', label='I')
plt.plot(t, R, 'k', label='R')
plt.legend()
plt.xlabel('Time (days)')
plt.ylabel('Number')
plt.show()


\end{verbatim}
\end{minipage}

\vspace{5mm}

\noindent
Does the maximum number of infected people agree with the value predicted
by \eqref{eq:Imax}?

\vspace{3mm}

\noindent
In section~\ref{sec:fpi} we found the value of $q$ needed to obtain $I_\mathrm{max} = 10^5$.
Set $r = aq$ for this value of $q$ and re-run the code; do we get $I_\mathrm{max} \approx 10^5$?
(You might need to run for more days.)


\section{Projectiles with air resistance: part~II}

We now have almost all the numerical methods we need to investigate projectiles
with air resistance. However, the time stepping methods we have seen apply to
first order ODEs whereas Newton's second law gives us second order ODEs. What
can we do?

\subsection{Time stepping higher order ODEs}
\label{sec:step_higher_systems}

Suppose we are given a second order ODE
\begin{equation}
\frac{d^2 y}{d t^2} = f(y, \dot{y}, t) .
\end{equation}
This can always be re-written as a pair of first order ODEs by introducing a
new variable $z = \dot{y}$, say:
\marginnote{
This pair of ODEs will need two initial conditions, one on $y$
and one on $z$. This makes it clear why the original second order ODE
also needs two initial conditions.
}
\begin{eqnarray}
\frac{d y}{d t} & = & z ; \nonumber \\
\frac{d z}{d t} & = & f(y, z, t) . \nonumber
\end{eqnarray}
We can apply any of the time stepping schemes we have met to this
pair of equations.

Clearly a similar trick can be used with higher order ODEs and with systems
of second order and higher order ODEs.


\subsection{Projectiles with air resistance}
\label{sec:air_resistance_examples}

To solve \eqref{eq:proj_air_res} numerically we can first split it into
horizontal and vertical components
\begin{eqnarray}
m \ddot{x} & = & \ \ \ \ \ \ \ \ \ - C |\boldsymbol{v}| \dot{x} , \nonumber \\
m \ddot{z} & = & -m g              - C |\boldsymbol{v}| \dot{z} , \nonumber
\end{eqnarray}
and then introduce new variables $u = \dot{x}$, $w = \dot{z}$ to obtain
a system of first order equations
\begin{eqnarray}
  \dot{x}  & = & u , \nonumber \\
  \dot{z}  & = & w , \nonumber \\
m \dot{u} & = & \ \ \ \ \ \ \ \ \ - C |\boldsymbol{v}| u , \nonumber \\
m \dot{w} & = & - m g             - C |\boldsymbol{v}| w , \nonumber
\end{eqnarray}
where $\boldsymbol{v} = (u,w)$ and $|\boldsymbol{v}| = (u^2 + w^2)^{1/2}$.
\marginnote{Note how the air resistance term couples the horizontal and vertical component
equations; we can't solve them separately as we did in section~\ref{sec:projectiles}}

Now we can use, for example, a Forward Euler method, as we did
in section~\ref{sec:step_systems} to obtain a numerical solution.
Since we have four equations we will need four initial conditions, for
example the $x$- and $z$-components of the initial position and velocity.

\ex{Write a Python code to solve these equations. You might want to
take the code of section~\ref{sec:step_systems} as a template.}

\vspace{3mm}

It is very important to test any code to make sure it does what it is supposed to.
For this problem we have closed form solutions in three cases (section~\ref{sec:air_res_I})
that we can compare our numerical solution against.

\vspace{3mm}

%\noindent
%Exercise. use arrays for x and v

\ex{Let us reconsider the problem from section~\ref{sec:projectiles},
  but now including air resistance. A ball is projected with speed
  $\sqrt{3gH}$ in a plane perpendicular to a vertical wall of height
  $H$ a distance $H$ away. Let us fix the magnitude of the drag
  coefficient $C$ by taking $C/m = 0.1/H$. Find the range of angles of
  projection that allow the ball to pass over the wall. We can no
  longer solve this problem analytically, but we can still answer the
  question by computing numerical solutions for different launch
  angles~$\alpha$.  It turns out that the solution is independent of
  the values given for $H$ and $g$, so we can choose to work in
  units\footnote{ A very useful technique in applied mathematics is
  \textit{non-dimensionalization}.  We won't say more about this here,
  but that is effectively what we have done.  } for which $H=1$,
  $g=1$.  It is a very good idea to check that your code gives the
  correct solution for the case of zero drag $C/m = 0$ first! Do you
  find that the range of angles to clear the wall increases or
  decreases when air resistance is included?}

\ex{We can also revisit the example from section~\ref{sec:envelope}.
  A ball is projected with speed $V$ in a plane perpendicular to a
  wall of height $H$ a distance $H$ away, with $C/m = 0.1/H$.  What is
  the {\it minimum} speed required for the ball to pass over the wall?
  Once again we can no longer solve the problem analytically, but we
  can obtain an answer using a trial-and-error approach with the code
  for the above example. (We will see another way to tackle this
  example in section~\ref{sec:adhoc}.)}
% 1.66*sqrt(gH) cf 1.5538*sqrt(gH) when C = 0.

\vspace{10mm}
